\subsection{One slope comparison over the translate group}\label{geo:translate}

Fix one good scale and one choice of the data of Formal input~\ref{red:design}.
Write \(\mathbb E a=|A_0|^{-1}\sum_{u\in A_0}a_u\). The translate \(u\)
and a word \(i\) at \(u\) are different indices. In a presentation with
several points above \(u\), select one lift for each \(u\) and use it
at every stage. We make this choice throughout. A probability mixture
of lifts gives the same proof if each lift's two rank measures are
multiplied by its conditional probability given \(u\).
The tied weights, the integers \(a_t\) and the centre are the ones on
these data. An identification of the points with translates in the
character group supplies no identification of their lattices.

The direct sum can be formed at the first edge. Retain the label \(u\)
in every factor and compose maps of the same word. For composable
layers \(C_{u,j}\), the identity
\[
 (\bigoplus_u C_{u,m})\cdots(\bigoplus_u C_{u,1})
     =\bigoplus_u(C_{u,m}\cdots C_{u,1})
\]
follows by evaluation on each summand. Over different coefficient fields
this is a tuple of maps over the product of those fields. The last two
kernels of the theta chain use the same intermediate mark \(Y\) in
each path. Treat their coupled sum as a single map whenever only that
sum descends to the quotient. If the fibres are identified, a diagonal
inclusion and a mean projection at the two \emph{ends} give the mean
of the composites. Averaging the edges separately can change the
answer even for integral isometries. On two rank-one blocks take \(+1\)
and \(-1\) at each of two successive edges in the same order. The
average composite is \(1\), whereas both average edges are zero.

For integral matrices in common bases, \(h^{-1}\sum_{u}B_u\),
\(h=|A_0|\), guarantees only the minimum entry valuation among the
summands minus \(v_w(h)\). On the present host \(h\) is a power of
two. Swapping the coordinates of \(\mathbb Z_2^2\) gives the
invariant projector \(e_1\mapsto(1/2,1/2)\). If a terminal integral
quotient \(q\) satisfies \(q\rho(u)=q\) for an integral group action
\(\rho\), then \(q\sum_u\rho(u)=hq\), so the division cancels there.
The construction below needs only the block maps. It divides real
degrees and traces by \(h\). It also keeps the weights inside the
local prime gains. For example, averaging \((2-w)_+\) at \(w=1,3\)
gives \(1/2\), while substitution of the average weight gives zero.

For word \(i\) write \(B_{ui}:S_{ui,K_u}\to R_{u,K_u}\) for the
rational descended column in the prescribed divisor shifts.
Here \(S_{ui},R_u\) are Hermitian lattices over a totally imaginary
field \(K_u\). Omit zero-rank words and put \(r_{ui}=\operatorname{rk}S_{ui}\),
\(r_u=\sum_i r_{ui}>0\), \(d_u=\operatorname{rk}R_u>0\). On the
tied data \(R_u\) is the dual target lattice of the column pairing. Slopes
are normalized as in Theorem~\ref{geo:hn}. A slope in word \(i\) has
mass \(1/(h r_u)\), and a target slope has mass \(1/(h d_u)\).
The ordinary slope of a direct sum over one field weights blocks by
rank. Repeating both sides of each block \(m_u\) times would require
both \(m_u r_u\) and \(m_u d_u\) to be constant to recover these
two means. The ratio \(r_u/d_u\) need not be constant. We use the
two normalized traces explicitly.

\begin{lemma}[Cost of making the descended column integral]\label{trmean:graph}
For one rational map \(B:S_K\to R_K\) fix real archimedean
allowances \(h_v\) with normalized mean \(h_B\). Write
\(\langle a_v\rangle_\infty=(2/[K:\mathbb Q])\sum_{v\mid\infty}a_v\)
with one representative of each conjugate pair. Define
\[
\begin{aligned}
 b_w(B)&=\operatorname{length}_{\mathcal O_{K,w}}
              ((R_w+BS_w)/R_w),\\
 p_v(B;h_v)&=\sum_a\log\max(1,e^{-h_v}\sigma_a(B_v)^2),\\
 g(B)&=2\sum_{w\nmid\infty}\frac{\log Nw}{[K:\mathbb Q]}b_w(B)
                    +\langle p_v(B;h_v)\rangle_\infty ,
\end{aligned}
\]
where \(\sigma_a\) are singular values for the prescribed metrics
and \(Nw\) is the residue field cardinality. There is a source
\(S^\sharp\) with integral contractive inclusion in \(S\) such that
\(B:S^\sharp\to R\) is integral and has height at most \(h_B\)
when it is nonzero. Its mean slope satisfies
\(\operatorname{rk}(S)(\mu(S)-\mu(S^\sharp))=g(B)\).
For composable maps, each \(b_w\) is subadditive and each \(p_v\)
is subadditive with additive allowances.
\end{lemma}
\begin{proof}
Take \(S^\sharp=S\cap B^{-1}R\) at the finite places. The map \(B\)
identifies \(S_w/S_w^\sharp\) with \((R_w+BS_w)/R_w\).
In Smith coordinates of the actual shifted lattices,
\(b_w=\sum_a(-\nu_{w,a})_+\) over the nonzero elementary divisors.
Only finitely many places contribute. The matrix
\(\operatorname{diag}(\pi^{-a},\pi^a)\) for \(a>0\) has determinant
one and \(b_w=a\). At \(v\) give the source the Gram matrix
\(\max(I,e^{-h_v}B_v^*B_v)\) in original orthonormal coordinates,
using the spectral decomposition. It dominates \(I\) and
\(e^{-h_v}B_v^*B_v\), proving the two norm bounds. Its log determinant
is \(p_v\). The finite index changes the squared normalized
covolume logarithm by \(2\sum_w b_w\log Nw/[K:\mathbb Q]\), and
the metric contributes \(\langle p_v\rangle_\infty\). Subtracting
the two normalized degrees proves the slope identity. Both changes
are zero for the zero map.

For \(X\xrightarrow{B}Y\xrightarrow{A}Z\) at a finite place,
\((Z+AY+ABX)/(Z+AY)\) is a quotient of \((Y+BX)/Y\).
The inclusion \(Z+ABX\subset Z+AY+ABX\) now gives
\(b_w(AB)\le b_w(A)+b_w(B)\). At infinity,
\[
 p_v(B;h_v)=\log\max_{k\ge0}
                 \|\textstyle\bigwedge^k(e^{-h_v/2}B_v)\|^2 .
\]
The term \(k=0\) equals one. The identity follows by ordering singular
values. Submultiplicativity of each exterior norm proves the assertion.
Thus positive costs may be accumulated from the first layer, before any
HN comparison, using the same intermediate lattices and metrics on the
two sides of each composition. In particular, if
\(B_{ui}=C_{ui,m}\cdots C_{ui,1}\) and the local allowances add to
those of the column, then
\[
 \mathbb E\,\frac{\sum_i g(B_{ui})}{r_u}
 \le \sum_{j=1}^m\mathbb E\,\frac{\sum_i g(C_{ui,j})}{r_u}.
\]
Each layer on the right uses the \emph{initial} total source rank
of its translate even when intermediate ranks differ. A coupled
terminal map counts as one factor in this formula.
\end{proof}

At one chosen infinite place in each block, suppose the original
source is \(V_{ui}\), with frame maps \(A_{ui}:V_{ui}\to\mathcal H_u\),
\(n_u=\dim\mathcal H_u\). The target identifies isometrically with
a subspace of dimension \(d_u\). Orthogonal projection of \(A_{ui}\)
to it gives \(B_{ui}\) up to a nonzero scalar, as in
Theorem~\ref{geo:hn}. Put \(Q_u=\sum_i A_{ui}A_{ui}^*\) and define
\[
\begin{aligned}
 e&=\mathbb E\,\frac{\sum_i\|A_{ui}^*A_{ui}-I\|_{\rm HS}^2}{r_u},
 &v&=\mathbb E\,\frac{\|(n_u/r_u)Q_u-I\|_{\rm HS}^2}{n_u},\\
 \delta&=\mathbb E(1-d_u/n_u),
 &\theta&=\min(1,\delta+\sqrt e+\sqrt v).
\end{aligned}
\]
Write \(h_{ui}\) for the mean archimedean allowance of \(B_{ui}\).
Let \(s_u=\sum_i r_{ui}(\mu(S_{ui})-h_{ui})/r_u\) and
\(\mu_u=\mu(R_u)\). Choose \(u\) uniformly, then choose a source
slope uniformly by rank in that block. Let \(S\) be its value
minus \(h_{ui}+s_u\). Let \(S'\) use the repaired source with the
\emph{same} centring constants. For the target let \(T\) be a
rank-uniform target slope minus \(s_u\). These specify three finite
probability laws. For such a law \(L\), \(L^\uparrow(p)\) denotes
its increasing quantile, so integrating on \([0,1]\) gives its mean.

\begin{theorem}[A single normalized trace comparison]\label{trmean:trace}
With these data, put
\(\mathcal J=\mathbb E\big(\sum_i g(B_{ui})/r_u\big)\). Then
\begin{equation}\label{trmean:gap}
 \mathbb E(s_u-\mu_u)\ \le\ \mathcal F,\qquad
 \mathcal F=\mathcal J+\int_{1-\theta}^1 S^\uparrow(p)\,dp
                         -\int_0^\theta T^\uparrow(p)\,dp .
\end{equation}
The two integrals use only the highest source fraction and the lowest
target fraction, in their original metrics. All ranks and frame
defects in this comparison may vary with the translate.
\end{theorem}
\begin{proof}
The HN cutoff of an integral map of height \(\le h\) sends source
slopes \(\ge a+h\) into the target cutoff at \(a\). Indeed, a nonzero
map from this source part into the quotient by that target part
would have an image which is a quotient of the former and whose
saturation is a sublattice of the latter. The quotient with quotient
metric has mean slope at least \(a+h\). The saturated image with
induced metric has mean slope \(<a\). The determinant of the induced
rank-\(k\) map lowers degree by at most \(kh\), since it is integral
and its archimedean squared norm is bounded by the \(k\)-th power of
the squared operator norm. This is a contradiction. This argument
uses the defining extremal properties of the HN filtration with its
subspace and quotient metrics. For the slope inequality and its
filtration form see the proof of Proposition~2.2 in
\href{https://arxiv.org/pdf/0803.2583}{H.~Chen, \emph{Positive degree and arithmetic bigness},
\S2.5}. Applied to \(S^\sharp\hookrightarrow S\)
at height zero it gives \(S'^\uparrow\le S^\uparrow\).
Lemma~\ref{trmean:graph} gives the three means \(0,-\mathcal J\),
and \(\mathbb E(\mu_u-s_u)\), respectively.

Fix a real threshold \(t\). Let \(U_{ui}\) be the subspace of the
repaired source above or at \(s_u+t+h_{ui}\), viewed in the original
complex vector space. The change of metric preserves its dimension
and its image as a subspace. Put \(q=\Pr(S'<t)\), let \(P_u\) project
orthogonally onto \((\sum_i A_{ui}U_{ui})^\perp\), and put
\(z=\mathbb E(\operatorname{tr}P_u/n_u)\). Then
\[
 z-\sqrt{zv}\ \le\ \mathbb E\,\frac{\operatorname{tr}(P_uQ_u)}{r_u}
                     \ \le\ q+\sqrt{qe}.
\]
For the left inequality pair \(P_u\) with \((n_u/r_u)Q_u-I\) in
the normalized Hilbert--Schmidt inner product, then apply
Cauchy--Schwarz over the blocks. For the right inequality
\(P_u A_{ui}\) vanishes on \(U_{ui}\). Sum the squared norms over
an orthonormal basis of its complement in the original source metric.
Pair its projection with \(A_{ui}^*A_{ui}-I\) and use the same
inequality with weights \(1/(h r_u)\). The mean relative dimension
of these complements is exactly \(q\).

If the relative codimension of the projected images in the target
is \(z_{0u}\), projection induces a surjection of the ambient
quotient onto the target quotient. Hence
\((d_u/n_u)z_{0u}\le\operatorname{tr}P_u/n_u\), so
\(\mathbb E z_{0u}\le z+\delta\). By the cutoff property, all projected
images lie in the target above or at \(s_u+t\).
Since \(z,q\le1\), we have
\(\Pr(T<t)\le\Pr(S'<t)+\theta\), including the case \(\theta=1\).
For finite increasing quantiles this implies
\(S'^\uparrow(p)\le T^\uparrow(p+\theta)\) for almost every
\(0<p<1-\theta\). Indeed, a strict violation yields a threshold
between the two values that contradicts the preceding count.
Integrate on that interval and use \(S'^\uparrow\le S^\uparrow\)
on the upper fraction. Substituting the three means gives
\eqref{trmean:gap}, also at the endpoints \(\theta=0,1\).
\end{proof}

A slope tail is essential. Take one block with \(D-1\) orthogonal unit
columns into the first \(D-1\) coordinates of a rank-\(D\) target,
with \(D\ge2\). All sources have slope zero on a common shifted scale,
and the target slopes are \(0,\ldots,0,-L\), obtained by changing
only the last line metric. Here \(\mathcal J=e=\delta=0\) and
\(v=1/(D-1)\), independent of \(L>0\). The mean gap is \(L/D\).
Consequently even small trace defects leave a quantitative tail
estimate to prove in the actual arithmetic family.

\begin{corollary}[Comparison with the error budget of the chain]\label{trmean:budget}
Suppose \(M_u=\operatorname{tr}Q_u\ge r_u/2\) and
\(d_u/n_u\ge1-\delta_*>0\) for the chosen blocks. Set
\[
 \eta=\mathbb E\left(\frac{n_u\operatorname{tr}(Q_u^2)}{M_u^2}-1\right),
 \qquad U=\mathbb E\max_{i,a}(\lambda_{ui,a}-h_{ui}-s_u),
\]
where \(\lambda_{ui,a}\) run over the original source slopes. For
\(H>0\) put \(T_{\rm out}=\int_0^1(-H-T^\uparrow(p))_+\,dp\).
Then \(U\ge0\) and
\begin{equation}\label{trmean:budget-eq}
 \mathbb E(s_u-\mu_u)\le T_{\rm out}+
 \frac{(\eta+4\sqrt e)H+4\mathcal J+4U\log(1+H/U)}{1-\delta_*}.
\end{equation}
The last product in the numerator is zero when \(U=0\). Integral
columns satisfying the prescribed global height bounds admit allowances
with \(\mathcal J=0\).
\end{corollary}
\begin{proof}
In the preceding trace proof normalize \(Q_u\) by \(M_u/n_u\). The
mean squared normalized defect is exactly \(\eta\), by expansion of
the square. The same projection test gives
\[
 z-\sqrt{z\eta}\le\mathbb E\,\operatorname{tr}(P_uQ_u)/M_u
                       \le 2(q+\sqrt{qe}).
\]
Here the upper bound uses \(M_u\ge r_u/2\) and positivity.
Since \(2\sqrt{z\eta}\le z+\eta\), the target fraction below the
threshold is at most \((\eta+4q+4\sqrt e)/(1-\delta_*)\).
At threshold \(-t\) the discarded source fraction is
\(q=\Pr(S'<-t)\). Its integral on \(0<t<H\) exceeds the corresponding
integral for \(S\) by at most \(\mathcal J\). Indeed the integrals
give the expectations of \(\min(H,(-S')_+)\) and \(\min(H,(-S)_+)\).
This function of the slope is decreasing and 1-Lipschitz, while the
ordered quantiles decrease by total integral \(\mathcal J\).

Write \(U_u=\max_{i,a}(\lambda_{ui,a}-h_{ui}-s_u)\). The centred
source in block \(u\) has mean zero and maximum \(U_u\ge0\).
For \(t>0\) its fraction below \(-t\) is at most \(U_u/(U_u+t)\).
The function \(a\mapsto a/(a+t)\) is concave on \(a\ge0\), so
\(\Pr(S<-t)\le U/(U+t)\). Integrating the three terms of the
target fraction on \([0,H]\) and adding \(T_{\rm out}\) bounds
the negative part of \(T\), which dominates \(-\mathbb E T\).
This proves \eqref{trmean:budget-eq}. For a nonzero integral column
of height \(\le h_{ui}\), choose \(h_v=\log\|B_v\|^2\) plus the
nonnegative constant making its normalized mean \(h_{ui}\).
Both terms of \(g(B)\) vanish. The zero column costs zero for any
allowances with the prescribed mean.
\end{proof}
